% concatenated from main.tex
\documentclass[11pt]{article}

\usepackage[margin=1.1in]{geometry}
\usepackage{amsmath,amssymb,amsthm}
\usepackage{hyperref}
\usepackage{comment}
\usepackage{graphicx} % Required for inserting images
\usepackage{listings}
\lstset{
basicstyle=\small\ttfamily,
columns=flexible,
breaklines=true
}

\newtheorem{theorem}{Theorem}
\newtheorem{lemma}{Lemma}

\DeclareMathOperator{\Sp}{Sp}
\newcommand{\ZZ}{\mathbb{Z}}
\newcommand{\QQ}{\mathbb{Q}}
\newcommand{\Gammaab}{\Gamma(\alpha,\beta)}

\title{Certifying Arithmeticity for Two Degree-Six Symplectic Hypergeometric Monodromy Groups}

\author{
  J. Maxwell Riestenberg\thanks{Max Planck Institute for Mathematics in the Sciences. Email:
		\href{mailto:riestenberg@mis.mpg.de} {\nolinkurl {riestenberg@mis.mpg.de}}.}
\and
  Diaaeldin Taha\thanks{Max Planck Institute for Mathematics in the Sciences. Email:
		\href{mailto:taha@mis.mpg.de} {\nolinkurl {taha@mis.mpg.de}}.} 
\and
  Steve Trettel\thanks{University of San Francisco. Email:
		\href{mailto:strettel@usfca.edu} {\nolinkurl {strettel@usfca.edu}}.} 
\and
  Anna Wienhard\thanks{Max Planck Institute for Mathematics in the Sciences. Email:
		\href{mailto:wienhard@mis.mpg.de} {\nolinkurl {wienhard@mis.mpg.de}}.} 
        \thanks{Center for Scalable Data Analytics and Artificial Intelligence, Leipzig, Germany}
}

\date{}

\begin{document}

\maketitle

\begin{abstract}
We prove arithmeticity for two degree-six symplectic hypergeometric monodromy groups, called C-47, and C-55 in the paper  \cite{BajpaiDonaNitsche2025Thin} by Bajpai-Dona-Nitsche. This settles two of the three remaining cases, whose classification was left open by \cite{BajpaiDonaNitsche2025Thin}. The arithmeticity certificates were found with AlphaEvolve and then independently verified with exact matrix arithmetic over $\QQ$ using a computer.
We include illustrations of limit sets of several degree-six symplectic hypergeometric monodromy groups. Based on these illustration we conjecture C-32 to be thin. 
\end{abstract}

\section{Introduction}
It is easy to give a subgroup of a matrix group such as $\mathrm{GL}(n,\mathbb{R})$, just write down some invertible matrices and consider the group $\Gamma$ they generate. When you define a group this way, it is however quite difficult in general to obtain information about the group $\Gamma$. It is discrete? And if yes, how big is it? What is its Zariski closure? Is it a lattice? 

One way to ensure discreteness is to pick matrices as generators which lie in an arithmetic lattice itself, for example in the group of integer matrices $\mathrm{SL}(n,\mathbb{Z})$ inside of $\mathrm{SL}(n,\mathbb{R})$. Then the main question becomes whether $\Gamma$ is of finite index in $\mathrm{SL}(n,\mathbb{Z})$ and thus itself a lattice, or of infinite index and thus a much smaller discrete subgroup. 
Such discrete subgroups of arithmetic lattices in $\mathrm{SL}(n,\mathbb{R})$ or more general semisimple matrix groups, which are Zariski dense, but of infinite index are called ``thin groups''. 

The names ``thin groups'' was introduced by Peter Sarnak who promoted their more than twenty years ago, motivated by the affine sieve methods developed by Bourgain-Gamburd-Sarnak, which allows to draw number theoretic conclusion for such groups, which were out of reach before. 

Particular nice families of discrete subgroups in arithmetic lattices arise from monodromy groups, see for example the discussion in Section~3.5. in \cite{Sarnak2014NotesThin}. They are often given by explicit generating matrices, and arise in interesting families. 

A particular nice family, which received a lot of attention over the past 20 years are monodromy representations of the classical hypergeometric equation.

Their monodromy representations are representations of the fundamental group of $\mathbb{C}\mathbb{P}^1 \backslash \{ 0, 1, \infty\}$ into $\mathrm{GL}_n(\mathbb{R})$, which are generated by the local monodromies around $0$, $1$, and $\infty$ given by matrices $A$, $B$, $C = A^{-1}B$. See 
\cite{Levelt1961HypergeometricFunctions,BeukersHeckman1989hypergeometric} for details. 

A special interest have been in such hypergeometric monodromy groups, which preserve a symplectic from $\Omega$ and are Zariski dense in $\mathrm{Sp}_{\Omega}(n,\mathbb{R})$.

The symplectic degree-four cases, which are connected to Calabi--Yau threefolds became
an early test family.  Brav--Thomas proved thinness for seven of them
\cite{BravThomas2014Thin}, while arithmeticity results of
Singh--Venkataramana and Singh settled the complementary cases
\cite{SinghVenkataramana2014Arithmeticity,Singh2015Four}.  In degree six,
Bajpai--Dona--Singh--Singh studied the corresponding symplectic
hypergeometric groups, motivated in part by a question of Katz about whether
the maximally unipotent degree-six family follows a similar pattern
\cite{BajpaiDonaSinghShashank2021degreesix}.  Bajpai--Dona--Nitsche later
settled many further degree-six cases, including many thin cases, by
computer-assisted ping-pong \cite{BajpaiDonaNitsche2025Thin}.  In their work \cite[Table 3]{BajpaiDonaNitsche2025Thin} they 
list the three remaining degree-six symplectic cases, C-32, C-47, and C-55, for
which neither arithmeticity nor thinness was known.


\begin{comment}
The arithmetic--thin dichotomy for monodromy groups is part of a broader
program popularized by Sarnak to understand Zariski-dense subgroups of
arithmetic groups that may nevertheless have infinite index
\cite{Sarnak2014NotesThin}.  In this setting, a subgroup can be as large as
possible algebraically, in the sense that its Zariski closure is the ambient
algebraic group, while still being small arithmetically.  Hypergeometric
monodromy groups provide a particularly concrete testing ground for this
problem: by Levelt's theorem they are generated by explicit companion
matrices, while the work of Beukers--Heckman determines their possible
Zariski closures
\cite{Levelt1961HypergeometricFunctions,BeukersHeckman1989hypergeometric}.

The arithmetic--thin problem for hypergeometric groups therefore becomes an
explicit finite-index problem.  Given an integral hypergeometric monodromy
group $\Gamma$ whose Zariski closure is an arithmetic group $G$, one asks
whether $\Gamma$ has finite index in $G(\ZZ)$, in which case it is arithmetic,
or infinite index, in which case it is thin.  This question is especially
tractable in the symplectic cases, where one can combine the
Beukers--Heckman classification with arithmeticity criteria for subgroups of
integral symplectic groups.

The symplectic degree-four cases connected to Calabi--Yau threefolds became
an early test family.  Brav--Thomas proved thinness for seven of them
\cite{BravThomas2014Thin}, while arithmeticity results of
Singh--Venkataramana and Singh settled the complementary cases
\cite{SinghVenkataramana2014Arithmeticity,Singh2015Four}.  In degree six,
Bajpai--Dona--Singh--Singh studied the corresponding symplectic
hypergeometric groups, motivated in part by a question of Katz about whether
the maximally unipotent degree-six family follows a similar pattern
\cite{BajpaiDonaSinghShashank2021degreesix}.  Bajpai--Dona--Nitsche later
settled many further degree-six cases, including many thin cases, by
computer-assisted ping-pong \cite{BajpaiDonaNitsche2025Thin}.  Their Table 3
lists three remaining degree-six symplectic cases, C-32, C-47, and C-55, for
which neither arithmeticity nor thinness was known.
\end{comment}

This note proves that two of these three remaining cases are arithmetic. We also include visualizations of the limit sets for all three cases, C-32, C-47, and C-55, which lead us to conjecture that C-32 is a thin group. However, we are unable to prove this at this point. 

To prove arithmeticity of the two groups 
C-47 and C-55, we use the arithmeticity criterion of Bajpai--Dona--Nitsche
\cite{BajpaiDonaNitsche2026Arithmetic} building on previous arithmeticity criteria of Venkataramana, Venkataramana-Singh
\cite{Venkataramana1987Zariskidensesubgroups, SinghVenkataramana2014Arithmeticity,Singh2015Four}. 
Their criterion reduces the problem to 
exhibiting explicit words in the  generators $A$ and $B$ of the hypergeometric monodropmy whose
conjugates of the standard rank-one unipotent give a  pair of commuting transvections.  The words
were found with AlphaEvolve \cite{NovikovEtAl2025AlphaEvolve,georgiev2025mathematical}.  Once the
words are fixed, the verification is deterministic and uses only exact matrix arithmetic over $\QQ$.

\section{Conventions and arithmeticity criterion}

Let $\alpha,\beta\in\QQ^6$ be parameter multisets.  The associated
cyclotomic polynomials are
\[
  f_\alpha(x)=\prod_{j=1}^6\bigl(x-\exp(2\pi i\alpha_j)\bigr),\qquad
  g_\beta(x)=\prod_{j=1}^6\bigl(x-\exp(2\pi i\beta_j)\bigr).
\]
For a degree-six monic polynomial
\[
  h(x)=x^{6}+c_5x^5+c_4x^4+c_3x^3+c_2x^2+c_1x+c_0,
\]
we use the companion matrix
\[
C(h)=
\begin{pmatrix}
0&0&0&0&0&-c_0\\
1&0&0&0&0&-c_1\\
0&1&0&0&0&-c_2\\
0&0&1&0&0&-c_3\\
0&0&0&1&0&-c_4\\
0&0&0&0&1&-c_5
\end{pmatrix}.
\]

The two groups considered in this note are the hypergeometric monodromy groups
labelled C-47 and C-55 in \cite[Table 3]{BajpaiDonaNitsche2025Thin}. We recall the definitions of these two groups. For
C-47, the parameter multisets are
\[
  \alpha_{47}=(0,0,1/5,2/5,3/5,4/5),\qquad
  \beta_{47}=(1/2,1/2,1/3,1/3,2/3,2/3),
\]
and the corresponding polynomials are
\[
  f_{47}=f_{\alpha_{47}}=x^6-x^5-x+1,\qquad
  g_{47}=g_{\beta_{47}}=x^6+4x^5+8x^4+10x^3+8x^2+4x+1.
\]
For C-55, the parameter multisets are
\[
  \alpha_{55}=(0,0,1/8,3/8,5/8,7/8),\qquad
  \beta_{55}=(1/2,1/2,1/12,5/12,7/12,11/12),
\]
and the corresponding polynomials are
\[
  f_{55}=f_{\alpha_{55}}=x^6-2x^5+x^4+x^2-2x+1,\qquad
  g_{55}=g_{\beta_{55}}=x^6+2x^5-2x^3+2x+1.
\]
The corresponding monodromy groups are
\[
  \Gamma_{47}=\langle C(f_{47}),C(g_{47})\rangle,\qquad
  \Gamma_{55}=\langle C(f_{55}),C(g_{55})\rangle.
\]
By the Beukers--Heckman classification
\cite{BeukersHeckman1989hypergeometric}, these primitive self-reciprocal
cyclotomic hypergeometric groups have Zariski closure equal to the
corresponding symplectic group.  Thus, after fixing the preserved integral
symplectic form $\Omega$ in each case, $\Gamma_{47}$ and $\Gamma_{55}$ are
Zariski-dense in $\Sp_\Omega(\mathbb{Z})$.

In either case, write
\[
  A=C(f),\qquad B=C(g),
\]
where $(f,g)$ is the corresponding pair of polynomials.  Also write
$a=A^{-1}$ and $b=B^{-1}$.  Words are read left to right as actions on column
vectors.  Thus, if $w=\ell_1\ell_2\cdots\ell_k$ with
$\ell_i\in\{A,B,a,b\}$, then the corresponding matrix is
\[
  M(w)=M(\ell_k)\cdots M(\ell_2)M(\ell_1),
\]
where $M(A)=A$, $M(B)=B$, $M(a)=A^{-1}$, and $M(b)=B^{-1}$.

The arithmeticity proofs for C-47 and C-55 use the following criterion of Bajpai--Dona--Nitsche
\cite[Lemma 1]{BajpaiDonaNitsche2026Arithmetic}, which builds on work of
Singh--Venkataramana \cite{SinghVenkataramana2014Arithmeticity} and
Venkataramana \cite{Venkataramana1987Zariskidensesubgroups}.

\begin{lemma}[Bajpai--Dona--Nitsche]\label{lem:bdn}
Let $n\geq 2$, let $\Omega$ be a nondegenerate symplectic form on $\QQ^{2n}$
which is integral on $\ZZ^{2n}$, and let
$\Gamma<\Sp_{\Omega}(\ZZ)$ be Zariski-dense. Then $\Gamma$ has finite index
in $\Sp_{\Omega}(\ZZ)$ if and only if $\Gamma$ contains two transvections
\[
  X_i=1+\lambda_i x_i\Omega(x_i,\cdot),\qquad
  \lambda_i\in\QQ^\times,\quad i=1,2,
\]
whose directions $x_1,x_2$ are linearly independent and
$\Omega$-orthogonal, i.e. $\Omega(x_1,x_2)=0$.
\end{lemma}

For the groups considered here, put
\[
  T=A^{-1}B.
\]
This is the standard rank-one unipotent element associated with the local
monodromy at $1$.  For a rank-one unipotent $U$, the one-dimensional subspace
$\operatorname{im}(U-I)$ is its transvection direction.  Since conjugates of
transvections are transvections, any word $\gamma$ in $A^{\pm1},B^{\pm1}$
gives another transvection $\gamma T\gamma^{-1}$.  Therefore, once
preservation of the relevant integral symplectic form has been checked,
Lemma \ref{lem:bdn} reduces arithmeticity to finding a word $\gamma$ such
that the directions of $T$ and $\gamma T\gamma^{-1}$ are linearly independent
and $\Omega$-orthogonal.

\section{C-47 and C-55 are arithmetic}

We now prove that C-47 and C-55 are arithmetic by verifying explicit
certificates found with AlphaEvolve.  Once the witness words are fixed, the
verification is deterministic and uses only exact matrix arithmetic over $\QQ$.

\begin{theorem}\label{thm:main}
The hypergeometric monodromy groups labelled C-47 and C-55 in
\cite[Table 3]{BajpaiDonaNitsche2025Thin} have finite index in their ambient
integral symplectic groups.  In particular, both groups are arithmetic.
\end{theorem}

\begin{proof}
First consider C-47.  Set
\[
  A=C(f_{47}),\qquad B=C(g_{47}),\qquad T=A^{-1}B.
\]
Let $w_{47}$ be the following freely reduced word of length $93$:
\begin{lstlisting}
bbbbbaabbbbbbAAbbbbbbaabbbbbbAAbbbbbbAAbbbbbbABaBaaBBBBBBAAAbAbaaBaBaBaBaBAAAAAAABaBaB
BaBabaa
\end{lstlisting}
and let $\gamma=M(w_{47})$ be the corresponding group element.

A direct computation over $\ZZ$ gives
\[
  \operatorname{im}(T-I)=\QQ x_1,\qquad
  \operatorname{im}(\gamma T\gamma^{-1}-I)=\QQ x_2,
\]
where
\[
  x_1=(-5,-8,-10,-8,-5,0),
\]
and
\[
  x_2=(491566906334,537748595482,224774947812,
       73905511690,-18977654566,0).
\]
The vectors $x_1$ and $x_2$ are linearly independent.
The integral
symplectic form preserved by $A$ and $B$ is
\[
\Omega_{47}=\begin{pmatrix}
0&29&-50&51&-28&1\\
-29&0&29&-50&51&-28\\
50&-29&0&29&-50&51\\
-51&50&-29&0&29&-50\\
28&-51&50&-29&0&29\\
-1&28&-51&50&-29&0
\end{pmatrix}.
\]
Indeed,
\[
  A^t\Omega_{47}A=\Omega_{47},\qquad
  B^t\Omega_{47}B=\Omega_{47},
\]
and
\[
  \det(\Omega_{47})=1679616\neq 0.
\]
Moreover,
\[
  \Omega_{47}(x_1,x_2)=0.
\]
The same exact computation gives that $T$ and $\gamma T\gamma^{-1}$ are
rank-one unipotents and that
\[
  T \, \gamma T \gamma^{-1} = \gamma T \gamma^{-1} \, T . %[T,\gamma T\gamma^{-1}]=I.
\]
Thus $T$ and $\gamma T\gamma^{-1}$ are commuting transvections whose
directions are linearly independent and $\Omega_{47}$-orthogonal.

Now consider C-55.  Set
\[
  A=C(f_{55}),\qquad B=C(g_{55}),\qquad T=A^{-1}B.
\]
Let $w_{55}$ be the following freely reduced word of length $49$:
\begin{lstlisting}
baaaabaaaabaaaabaaaabaaaabaaaabaaaaaBABaBABAAbaaB
\end{lstlisting}
and let $\gamma=M(w_{55})$ be the corresponding group element.

A direct computation over $\ZZ$ gives
\[
  \operatorname{im}(T-I)=\QQ x_1,\qquad
  \operatorname{im}(\gamma T\gamma^{-1}-I)=\QQ x_2,
\]
where
\[
  x_1=(-4,1,2,1,-4,0),
\]
and
\[
  x_2=(40999920,-275447328,-132048384,
       236325024,314749968,0).
\]
The vectors $x_1$ and $x_2$ are linearly independent. 
The integral
symplectic form preserved by $A$ and $B$ is
\[
\Omega_{55}=\begin{pmatrix}
0&1&6&3&4&5\\
-1&0&1&6&3&4\\
-6&-1&0&1&6&3\\
-3&-6&-1&0&1&6\\
-4&-3&-6&-1&0&1\\
-5&-4&-3&-6&-1&0
\end{pmatrix}.
\]
Indeed,
\[
  A^t\Omega_{55}A=\Omega_{55},\qquad
  B^t\Omega_{55}B=\Omega_{55},
\]
and
\[
  \det(\Omega_{55})=4096\neq 0.
\]
Moreover,
\[
  \Omega_{55}(x_1,x_2)=0.
\]
The same exact computation gives that $T$ and $\gamma T\gamma^{-1}$ are
rank-one unipotents and that
\[
  T \, \gamma T \gamma^{-1} = \gamma T \gamma^{-1} \, T . %[T,\gamma T\gamma^{-1}]=I.
\]
Thus $T$ and $\gamma T\gamma^{-1}$ are commuting transvections whose
directions are linearly independent and $\Omega_{55}$-orthogonal.

In both cases, the displayed form is integral and nondegenerate, and the
identities above show that the corresponding group
$\Gamma=\langle A,B\rangle$ lies in $\Sp_\Omega(\ZZ)$.  The Zariski-density
hypothesis in Lemma \ref{lem:bdn} is supplied by the Beukers--Heckman
classification, as recalled above.  Lemma \ref{lem:bdn} therefore applies to
C-47 and C-55, and shows that each group has finite index in its ambient
integral symplectic group.
\end{proof}

The accompanying
scripts\footnote{Available at \url{https://github.com/ditahd/Sp-6-Certificates}} reproduce the exact computations used in the proof. Precisely, they each reconstruct $A$, $B$, the
symplectic form, the witness word, and its inverse.
They verify over $\ZZ$ and $\QQ$ that $A$ and $B$ preserve the displayed form,
that $T$ and $\gamma T\gamma^{-1}$ are rank-one unipotents, that the two
elements commute, and that their directions are linearly independent and
$\Omega$-orthogonal.





\section{Limit set visualizations}\label{sec:limitset}


The proximal limit set $\Lambda(\Gammaab)$ of $\Gammaab = \langle A, B \rangle$
is the closure in $\mathbb{RP}^5$ of the attracting fixed lines of all
loxodromic elements of $\Gammaab$, and has visibly different qualitative
signatures in the two regimes: a proper compact subset (often fractal) in
the thin case versus the entire $\mathbb{RP}^5$ in the arithmetic case.  We
accompany the algebraic results with visualizations of these limit sets;
comparison with the limit set for C-32 forms the basis for its conjectured
thinness.  






\paragraph{How we compute it.}

We approximate $\Lambda(\Gammaab)$ by a finite \emph{partial
orbit}: starting from one attracting fixed line $\xi_+ \in \Lambda(\Gammaab)$ of a loxodromic $\eta$, we
apply to it the set of freely reduced words of length at most $N$ in
$\{A^{\pm 1}, B^{\pm 1}\}$.  
For the loxodromic element we take $\eta = TBT$, with $T = A^{-1}B$ as previously.  In all five examples $\eta$ has a real, simple,
dominant top eigenvalue $\lambda_1$ with spectral gap $|\lambda_1/\lambda_2| \geq 10$,
so the iteration
\[
  v_{k+1} \;=\; \frac{\eta\, v_k}{\|\eta\, v_k\|}
\]
starting from a generic seed converges projectively to $\xi_+$ to working
precision in under $30$ steps.

We walk through the set of words of length $\leq N$ inductively and compute
the corresponding vectors in the orbit $\mathcal{O}_N$ at each step. For a word
$w = w'\, s$ the representative $w \cdot \xi_+ \in \mathbb{R}^6$ is
obtained from the parent's representative $w' \cdot \xi_+$ by a single
matrix-vector product with the generator $s$, then rescaled to unit norm
(preserving the projective class). 
In the images, we color points by the final generator applied with $a=$red, $A=$yellow, and $b=$green, $B=$blue.
Our visualizations use words of length up to 21, yielding orbits of approximately 20 billion points.  



\paragraph{Projection $\mathbb{RP}^5 \dashrightarrow \mathbb{R}^3$.}
Fix $\ell \in (\mathbb{R}^6)^*$ and project onto the affine chart
$H_\ell = \{v : \ell(v) = 1\}$, discarding representatives with
$|\ell(v)| < 10^{-3}$ near infinity in the patch.  Inside $H_\ell$ we compute the principal components
of the centered cloud $\{v/\ell(v) : v \in \mathcal{O}_N\}$ and plot its
projection onto the top three.  Using principal components ensures that the
three plotted axes capture the directions of greatest variance in the
cloud, so that an orbit that genuinely spreads in more than three
dimensions is not artificially flattened by the projection.




\paragraph{The figures.}
We compare three groups of pictures.  Figure~\ref{fig:calibration} fixes
the visual templates of arithmetic and thin, showing the partial orbits
for two cases classified by \cite{BajpaiDonaNitsche2025Thin}: in their enumeration, A-1 (thin) and A-15 (arithmetic).  The proximal limit set of A-1 lies near a thin curve whereas even just the first million points for A-15 spread out in many directions.

Figure~\ref{fig:arithmetic} shows the
partial orbits for C-47 and C-55, the two cases proved arithmetic in
Theorem~\ref{thm:main}; both clouds can also be seen to be more complex than that of A-1.  Finally, Figure~\ref{fig:c32} shows the
partial orbit for C-32 appearing as a curve,
qualitatively closer to A-1 than to A-15 or C-47, C-55 images, suggesting it may be thin.

\begin{figure}[t]
  \centering
  \includegraphics[width=0.49\linewidth]{img/a1.jpg}\hfill
  \includegraphics[width=0.49\linewidth]{img/a15.jpg}
  \caption{Partial limit sets for A-1 (left, thin)
    and A-15 (right, arithmetic) containing $\approx 20$ billion points.}
  \label{fig:calibration}
\end{figure}

\begin{figure}[t]
  \centering
  \includegraphics[width=0.49\linewidth]{img/c47.jpg}\hfill
  \includegraphics[width=0.49\linewidth]{img/c55.jpg}
  \caption{Partial limit sets for C-47 (left) and C-55
    (right) containing  $\approx 20$ billion points.}
  \label{fig:arithmetic}
\end{figure}

\begin{figure}[t]
  \centering
  \includegraphics[width=0.49\linewidth]{img/c32a.jpg}\hfill
   \includegraphics[width=0.49\linewidth]{img/c32b.jpg}
  \caption{Partial limit set for C-32, containing 20 billion points. The right is a deep deep zoom.}
  \label{fig:c32}
\end{figure}










\section*{Acknowledgements}

The witness words in this note were found with AlphaEvolve.  DT thanks the
AlphaEvolve team at Google DeepMind for access to the system and for
discussions about auditable AI-assisted mathematical discovery.
AW thanks the Hector Fellow Academy for support.


\clearpage

\bibliographystyle{amsalpha}
\bibliography{biblio}

\end{document}
